\input{../../../templates/es/preambulo.tex}
\setbeamertemplate{navigation symbols}{}
\setbeamertemplate{footline}[frame number]
\date{}
\begin{document}
\tituloleccion{04}{Cruza}
\begin{frame}{Pregunta de diseño}
\normalsize
Dados dos padres fijos, ¿qué descendientes puede alcanzar un operador? El caso de Planta Física usa 48 decisiones reales, cuatro edificios por doce meses.
\end{frame}
\begin{frame}{Qué se mide}
\normalsize
El alcance es el soporte de la distribución de hijos. Mido hijos distintos, genes nuevos, clones y fracción legal. Una muestra ilustra; enumerar demuestra el tamaño del soporte.
\end{frame}
\begin{frame}{Línea base}
\normalsize
\begin{center}\includegraphics[width=.84\textwidth,height=.53\textheight,keepaspectratio]{../figures/descendencia_01_dos_padres.png}\end{center}
{\small Los dos padres fijos dan una línea base comparable para cada operador.}
\end{frame}
\begin{frame}{Un punto}
\normalsize
Con seis genes distintos, cinco cortes interiores y dos orientaciones producen exactamente 10 hijos. Cada gen se copia de uno de los padres: no hay valores nuevos.
\end{frame}
\begin{frame}{Varios cortes y uniforme}
\normalsize
\begin{center}\includegraphics[width=.84\textwidth,height=.53\textheight,keepaspectratio]{../figures/descendencia_03_mas_cortes.png}\end{center}
{\small En esta implementación: un corte 10 hijos, dos cortes 12, tres cortes 8, uniforme 64.}
\end{frame}
\begin{frame}{Un detalle del algoritmo}
\normalsize
Los tres cortes se toman de \texttt{range(1,L-1)}: se excluye la última frontera interior. El resultado 8 depende de esa implementación; no es una ley general sobre más cortes.
\end{frame}
\begin{frame}{BLX-\(\alpha\)}
\normalsize
Si \(m=\min(a,b)\), \(M=\max(a,b)\) y \(d=M-m\), entonces \(x'\sim U[m-\alpha d,M+\alpha d]\). Para \(\alpha=.5\), salir del intervalo parental tiene probabilidad \(1/2\).
\end{frame}
\begin{frame}{Acotar no basta}
\normalsize
\begin{center}\includegraphics[width=.84\textwidth,height=.53\textheight,keepaspectratio]{../figures/descendencia_04_mezcla.png}\end{center}
{\small La mezcla calcula genes nuevos y puede salir del dominio. Acotar sólo repara límites por gen.}
\end{frame}
\begin{frame}{Guiada por aptitud}
\normalsize
Evaluar candidatos y conservar los que no empeoran reduce el deterioro observado. Cada intento interno consume una llamada al evaluador; devolver un padre sin cambio también es posible.
\end{frame}
\begin{frame}{Representación}
\normalsize
Un vector real admite mezcla aritmética. Una permutación debe visitar cada parada exactamente una vez. Un corte ordinario puede repetir paradas y omitir otras.
\end{frame}
\begin{frame}{Rutas ilegales}
\normalsize
\begin{center}\includegraphics[width=.84\textwidth,height=.53\textheight,keepaspectratio]{../figures/descendencia_06_la_division.png}\end{center}
{\small Corte y mezcla no construyen rutas legales en esta prueba. La cruza uniforme sólo recupera un padre legal.}
\end{frame}
\begin{frame}{Cruza ordenada}
\normalsize
OX1 copia un segmento y llena los huecos en el orden relativo del otro padre. Conserva pertenencia y unicidad en permutaciones; sobre vectores reales puede resultar inerte.
\end{frame}
\begin{frame}{Dos familias}
\normalsize
\begin{center}\includegraphics[width=.84\textwidth,height=.53\textheight,keepaspectratio]{../figures/descendencia_07_ordenada.png}\end{center}
{\small OX1 funciona para rutas; una cruza aritmética funciona para valores reales. Elegir la familia es parte del modelo.}
\end{frame}
\begin{frame}{Transferencia al reto}
\normalsize
Para las 48 proporciones del caso, una cruza real propone un hijo. El evaluador compartido comprueba límites por gen, cuota anual, mínimo potable y variación mensual. Acotar cada gen no demuestra factibilidad global.
\end{frame}
\begin{frame}{Ficha de diseño}
\normalsize
En equipos de tres: elegir familia, parámetro, reparación y medida de diversidad. El revisor explica el método; el verificador cuenta evaluaciones y comprueba factibilidad.
\end{frame}
\begin{frame}{Interpretación}
\normalsize
Fracción legal no equivale a aptitud. Comparar operadores con el mismo presupuesto, ocho semillas y una política de reparación común. Estos resultados son sintéticos, no parámetros confirmados por Planta Física.
\end{frame}
\end{document}
